\begin{abstract}
India's courts are buried under nearly fifty million pending cases. The scale alone makes a strong case for computational intervention, yet the legal AI tools available to Indian practitioners today are built in isolation from one another: retrieval engines that know nothing about prediction models, prediction models that operate independently of explainability modules, and explainability components that rarely feed their signals back into retrieval. In this work, we survey Indian legal NLP through five capability axes, namely \emph{precedent search}, \emph{case fact structuring}, \emph{argument graph construction}, \emph{judicial outcome simulation}, and \emph{counterfactual what-if analysis}. For each axis we catalogue representative methods, benchmark datasets, and evaluation practices. Our principal observation is blunt: no existing system unifies all five axes under a single coherent architecture, and at least two of them, argument graph construction and counterfactual reasoning, are for practical purposes absent from the Indian legal AI landscape. Whether this fragmentation is best read as an artefact of how the field grew up, axis by axis, or as a more fundamental modelling oversight, we will not pretend to have fully settled. What we do argue is that further progress on individual axes appears to have reached a point of diminishing returns, and that meaningful advancement now depends on a unifying pipeline that couples retrieval, structured reasoning, prediction, and interactive analysis.
\end{abstract}

%% ================================================================
%%  SECTION 1 — Introduction
%% ================================================================
\section{Introduction}

\IEEEPARstart{I}{ndia's} courts are drowning. Roughly fifty million cases sit pending across district courts, high courts, and the Supreme Court of India, and the backlog grows each year faster than the judiciary can clear it. This is not, we believe, a scheduling problem; it is a structural failure of access to justice that periodic administrative reform has been unable to fix. Litigants wait years, sometimes a full decade, for resolution, and public trust erodes quietly, case by case. In this setting, AI-assisted legal tools offer a way to augment, not replace, judicial reasoning at a scale that human capacity alone cannot match. Over one million Supreme Court and High Court judgments are now digitised~\cite{nigam-etal-2025-nyayaanumana}, providing an empirical base that few other common-law jurisdictions can rival. We believe these conditions make India one of the most consequential, and still underexplored, testbeds for legal AI research today.

The tools that do exist, however, talk past one another. A retrieval engine that surfaces relevant precedents has no awareness of the prediction model sitting one API call away. The prediction model, in turn, generates a disposition label without consulting how the retrieval module ranked its candidates. Explainability components, where they exist at all, are bolted on after the fact. This fragmentation, we suspect, does not match how judicial deliberation actually unfolds in practice — though we should be honest that the line between how judges reason and how judges describe their reasoning in written opinions is one we have not fully resolved, and the literature we draw on here largely reads the latter as evidence for the former. A judge does not first retrieve, then apply statute, then weigh facts; she does all three at once, circling back the moment a new precedent reshapes the factual picture. Recent empirical evidence aligns with this intuition: coupling retrieval with downstream reasoning outperforms strict pipeline decomposition~\cite{nigam-etal-2025-nyayarag}, and case-based reasoning integrated into retrieval-augmented generation produces answers that are measurably more grounded in legal sources~\cite{wiratunga-etal-2024-cbrrag}. The lesson, we take it, is that tight integration across capability boundaries is not optional but constitutive of legally faithful reasoning.

We should be upfront about what this paper is. It is a survey, not a system. Our contributions are twofold. First, we offer a structured review of five capability areas in Indian legal NLP, with attention to the methods, datasets, and evaluation practices that have shaped each. Second, we identify the missing integrated multi-capability pipeline as the field's central open problem, and we sketch what such a pipeline would need to do. The five axes are: (1)~\emph{Precedent Search}, encompassing sparse and dense retrieval of prior judgments; (2)~\emph{Case Fact Structuring}, including named entity recognition, rhetorical role labelling, and statutory section tagging; (3)~\emph{Argument Graph Construction}, covering the extraction of claims, evidence, and their relational structure; (4)~\emph{Judicial Outcome Simulation}, predicting case dispositions alongside natural-language explanations; and (5)~\emph{What-If Analysis}, enabling counterfactual hypothesis testing over case facts. Taken together, these axes span the full reasoning chain from information retrieval to interactive deliberation.

The remainder of this paper follows a straightforward path. Section~II sets up background on the Indian legal ecosystem and the specific difficulties it creates for NLP. Sections~III through VII work through each capability axis in turn. Section~VIII steps back to compare across axes, flagging integration gaps and recurring themes. Section~IX takes up evaluation methodologies. Section~X outlines open problems and future directions. Section~XI concludes.

\begin{table*}[t]
\centering
\caption{Five Capability Axes in Indian Legal AI: Coverage and Identified Gaps}
\label{tab:axes}
\begin{tabular}{>{\raggedright\arraybackslash}p{3.2cm} >{\raggedright\arraybackslash}p{6.8cm} >{\raggedright\arraybackslash}p{6.8cm}}
\toprule
\textbf{Capability Axis} & \textbf{Description} & \textbf{Gap in Existing Literature} \\
\midrule
Precedent Search &
Semantic retrieval of relevant prior judgments through sparse keyword-based methods, dense vector representations, and hybrid strategies applied to Indian case law corpora. &
Retrieval systems function as standalone components; their outputs are not consumed by downstream prediction or reasoning modules, limiting their utility in end-to-end legal analysis. \\
\addlinespace
Case Fact Structuring &
Extraction and classification of structured elements from judgment text, including named entity recognition, rhetorical role labelling of sentence-level functional segments, and Indian Penal Code section tagging. &
Structured representations produced by these models are not propagated to any subsequent reasoning stage, rendering them terminal outputs rather than intermediate features in a larger pipeline. \\
\addlinespace
Argument Graph Construction &
Identification and relational mapping of argumentative components, including claims, supporting evidence, cited statutes, and invoked precedents, into directed graph representations that capture the logical structure of judicial reasoning. &
This capability is virtually absent in Indian legal AI research; existing work on legal argument mining is confined almost exclusively to Western jurisdictions such as the European Court of Human Rights. \\
\addlinespace
Judicial Outcome Simulation &
Prediction of case dispositions (acceptance or dismissal) accompanied by natural language explanations that articulate the reasoning underlying the predicted outcome. &
Prediction models operate without access to retrieval-augmented context; generated explanations remain shallow and lack grounding in specific precedents or statutory provisions. \\
\addlinespace
What-If Analysis &
Counterfactual hypothesis testing that enables users to perturb specific case facts, statutory provisions, or precedent selections and observe the resulting shifts in predicted outcomes and explanatory rationales. &
Entirely unexplored in Indian legal NLP; no existing dataset, model, or evaluation framework addresses counterfactual reasoning over Indian case law. \\
\bottomrule
\end{tabular}
\end{table*}

%% ================================================================
%%  SECTION 2 — Background: The Indian Legal Ecosystem
%% ================================================================
\section{Background: The Indian Legal Ecosystem}

India inherited its common law tradition from British colonial rule, and that inheritance carries real weight in everyday practice: judicial decisions are precedent, not merely advisory. Judicial reasoning here rests on three pillars that a court must hold in mind at once, namely the established facts of the case, the applicable statutory provisions (principally the Indian Penal Code and the Code of Civil Procedure), and binding prior rulings invoked through the doctrine of \emph{stare decisis}. No competent judge consults these sequentially. A judgment forms at the intersection of all three: facts are read through the lens of statute, statutes are interpreted in light of prior rulings, and prior rulings are tested against the factual matrix of the case at hand. We contend that any computational framework which treats these pillars as independent inputs to a staged pipeline has, almost by construction, already compromised its fidelity to how law actually operates.

\IEEEPARstart{P}{rocessing} Indian legal text introduces a set of difficulties that researchers working on ECHR or U.S. federal case law rarely encounter. The most immediate is multilingualism. A single judgment may shift between English, Hindi, Urdu, and a regional language mid-paragraph, with no markup or translation cues, which is a serious obstacle for tokenisers trained on monolingual corpora. Formatting is another persistent pain point. There is no standard template across Indian courts; judgments from different benches, different decades, and different levels of the judicial hierarchy vary substantially in structure, citation conventions, and rhetorical organisation, and rule-based extraction tends to break almost immediately under this variation. A third difficulty is annotation scarcity. The LegalEval shared task made this concrete: even trained annotators disagreed substantially on rhetorical role labels for Indian legal text~\cite{modi-etal-2023-semeval}. General-purpose encoders such as base BERT are not adequate to the task either. InLegalBERT demonstrated that domain-adapted pre-training on Indian corpora is essentially non-negotiable for competitive performance~\cite{paul-etal-2023-inlegalbert}. We suspect the community still underestimates how much domain shift matters in this setting — or, more honestly, perhaps the better way to put it is that domain shift is widely acknowledged in the abstract but consistently underfunded in practice, and we are not yet sure which framing leads to the right intervention.

The good news is that resource creation has accelerated sharply over the past five years. ILDC offered the community its first large-scale benchmark, providing over thirty-five thousand Supreme Court cases annotated for judgment prediction and explanation~\cite{malik-etal-2021-ildc}. PredEx raised the bar by attaching expert-annotated explanations to specific textual spans within judgments~\cite{nigam-etal-2024-predex}. NyayaAnumana, the most recent and largest corpus to date, spans multiple court levels and ships with its own domain-specialised language model~\cite{nigam-etal-2025-nyayaanumana}. Together, these datasets have driven measurable gains on tasks such as binary judgment prediction and rhetorical role classification. However, every one of them targets a single capability axis in isolation. No existing work unifies retrieval, structuring, reasoning, prediction, and interactive analysis under one coherent architecture. That gap, the absence of integration rather than the absence of data, is what motivates the five-axis taxonomy we develop in the sections that follow.
